% ========================================================
% BÀI TẬP TOÁN 5 - KẾT NỐI TRI THỨC VỚI CUỘC SỐNG
% Bài 17: Thực hành và trải nghiệm với một số đơn vị đo đại lượng
% Định dạng: BTTeX Part III - Trắc nghiệm điền số ngắn
% ========================================================

% Câu 1
\begin{ex}
  Một lớp học có chiều dài khoảng $8\text{ m}$, chiều rộng khoảng $6\text{ m}$. Diện tích phòng học đó ước lượng là bao nhiêu mét vuông?
  \shortans{48}
  \loigiai{
    Diện tích ước lượng: $8 \times 6 = 48\text{ m}^2$.
  }
\end{ex}

% Câu 2
\begin{ex}
  Một bàn học sinh hình chữ nhật có kích thước $120\text{ cm} \times 50\text{ cm}$. Diện tích mặt bàn là bao nhiêu đề-xi-mét vuông?
  \shortans{60}
  \loigiai{
    $120\text{ cm} = 12\text{ dm}$, $50\text{ cm} = 5\text{ dm}$. Diện tích là $12 \times 5 = 60\text{ dm}^2$.
  }
\end{ex}

% Câu 3
\begin{ex}
  Một bảng viết lớp học dài $3\text{ m}$, rộng $1{,}2\text{ m}$. Diện tích mặt bảng là bao nhiêu phần mười mét vuông? (Điền tử số khi mẫu số là 10, tức diện tích nhân 10).
  \shortans{36}
  \loigiai{
    Diện tích là $3 \times 1{,}2 = 3{,}6\text{ m}^2 = \frac{36}{10}\text{ m}^2$. Số cần điền là 36.
  }
\end{ex}

% Câu 4
\begin{ex}
  Ước lượng chiều dài sân bóng đá mini là $40\text{ m}$, chiều rộng là $20\text{ m}$. Diện tích sân bóng đó là bao nhiêu mét vuông?
  \shortans{800}
  \loigiai{
    Diện tích sân bóng là $40 \times 20 = 800\text{ m}^2$.
  }
\end{ex}

% Câu 5
\begin{ex}
  Sân trường hình chữ nhật có chu vi $180\text{ m}$, chiều dài hơn chiều rộng $10\text{ m}$. Diện tích sân trường là bao nhiêu mét vuông?
  \shortans{2000}
  \loigiai{
    Nửa chu vi là $90\text{ m}$. Chiều dài là $(90 + 10) : 2 = 50\text{ m}$. Chiều rộng là $40\text{ m}$. Diện tích là $50 \times 40 = 2\,000\text{ m}^2$.
  }
\end{ex}

% Câu 6
\begin{ex}
  Một quyển sách giáo khoa Toán 5 có kích thước khoảng $19\text{ cm} \times 26{,}5\text{ cm}$. Chu vi bìa sách là bao nhiêu xăng-ti-mét?
  \shortans{91}
  \loigiai{
    Chu vi là $(19 + 26{,}5) \times 2 = 45{,}5 \times 2 = 91\text{ cm}$.
  }
\end{ex}

% Câu 7
\begin{ex}
  Một viên gạch lát vỉa hè hình lục giác đều được ghép từ các tam giác đều nhỏ, diện tích mỗi tam giác là $15\text{ cm}^2$. Diện tích viên gạch lục giác là bao nhiêu xăng-ti-mét vuông?
  \shortans{90}
  \loigiai{
    Lục giác đều gồm 6 tam giác đều: $15 \times 6 = 90\text{ cm}^2$.
  }
\end{ex}

% Câu 8
\begin{ex}
  Một bồn hoa hình tam giác có đáy dài $4\text{ m}$, chiều cao tương ứng là $3\text{ m}$. Diện tích bồn hoa là bao nhiêu mét vuông?
  \shortans{6}
  \loigiai{
    Diện tích bồn hoa là: $(4 \times 3) : 2 = 6\text{ m}^2$.
  }
\end{ex}

% Câu 9
\begin{ex}
  Để trồng cỏ cho sân chơi diện tích $150\text{ m}^2$, chi phí mỗi mét vuông là $40\,000$ đồng. Hỏi tổng chi phí trồng cỏ là bao nhiêu triệu đồng?
  \shortans{6}
  \loigiai{
    Tổng chi phí là: $150 \times 40\,000 = 6\,000\,000$ đồng = 6 triệu đồng.
  }
\end{ex}

% Câu 10
\begin{ex}
  Một bể bơi mini hình chữ nhật dài $15\text{ m}$, rộng $6\text{ m}$. Người ta rào xung quanh bể bơi cách mép bể $1\text{ m}$. Chiều dài hàng rào là bao nhiêu mét?
  \shortans{50}
  \loigiai{
    Chiều dài khu đất có rào là $15 + 2 = 17\text{ m}$. Chiều rộng là $6 + 2 = 8\text{ m}$. Chu vi hàng rào là $(17 + 8) \times 2 = 50\text{ m}$.
  }
\end{ex}

% Câu 11
\begin{ex}
  Một tấm pin mặt trời hình chữ nhật có kích thước $2\text{ m} \times 1\text{ m}$. Một hộ gia đình lắp $10$ tấm pin như thế. Tổng diện tích các tấm pin là bao nhiêu mét vuông?
  \shortans{20}
  \loigiai{
    Diện tích 1 tấm pin là $2 \times 1 = 2\text{ m}^2$. 10 tấm có diện tích là $2 \times 10 = 20\text{ m}^2$.
  }
\end{ex}

% Câu 12
\begin{ex}
  Ước lượng cân nặng của một học sinh lớp 5 khoảng $35\text{ kg}$. Tổng cân nặng của một nhóm $6$ bạn học sinh như thế khoảng bao nhiêu ki-lô-gam?
  \shortans{210}
  \loigiai{
    Tổng cân nặng là: $35 \times 6 = 210\text{ kg}$.
  }
\end{ex}

% Câu 13
\begin{ex}
  Một bước chân của bạn Nam dài khoảng $50\text{ cm}$. Nam đi dọc theo hành lang lớp học hết $80$ bước chân. Hành lang lớp học dài khoảng bao nhiêu mét?
  \shortans{40}
  \loigiai{
    Chiều dài hành lang là: $50 \times 80 = 4\,000\text{ cm} = 40\text{ m}$.
  }
\end{ex}

% Câu 14
\begin{ex}
  Một hộp sữa chua cân nặng $100\text{ g}$. Một lốc sữa chua có $4$ hộp. Hỏi $5$ lốc sữa chua như thế cân nặng bao nhiêu ki-lô-gam?
  \shortans{2}
  \loigiai{
    Tổng số hộp là $4 \times 5 = 20$ hộp. Khối lượng là $20 \times 100 = 2\,000\text{ g} = 2\text{ kg}$.
  }
\end{ex}

% Câu 15
\begin{ex}
  Một chai nước khoáng chứa $500\text{ ml}$ nước. Một thùng nước khoáng gồm $24$ chai. Thùng nước khoáng đó chứa bao nhiêu lít nước?
  \shortans{12}
  \loigiai{
    Thể tích là $500 \times 24 = 12\,000\text{ ml} = 12$ lít.
  }
\end{ex}

% Câu 16
\begin{ex}
  Một gia đình dùng hết $15\text{ m}^3$ nước trong một tháng. Biết $1\text{ m}^3 = 1\,000\text{ l}$. Gia đình đó đã dùng hết bao nhiêu nghìn lít nước?
  \shortans{15}
  \loigiai{
    $15\text{ m}^3 = 15\,000$ lít = 15 nghìn lít.
  }
\end{ex}

% Câu 17
\begin{ex}
  Mỗi ngày bạn Lan uống khoảng $2$ lít nước. Hỏi trong tháng Tư (có $30$ ngày), bạn Lan uống khoảng bao nhiêu lít nước?
  \shortans{60}
  \loigiai{
    Lượng nước uống là: $2 \times 30 = 60$ (lít).
  }
\end{ex}

% Câu 18
\begin{ex}
  Một chiếc xe chở $40$ thùng hàng, mỗi thùng nặng $25\text{ kg}$. Tổng khối lượng hàng hóa trên xe là bao nhiêu tấn?
  \shortans{1}
  \loigiai{
    Tổng khối lượng là: $40 \times 25 = 1\,000\text{ kg} = 1\text{ tấn}$.
  }
\end{ex}

% Câu 19
\begin{ex}
  Một khu đất cắm trại hình chữ nhật có chiều dài $60\text{ m}$, chiều rộng $40\text{ m}$. Người ta chia đều thành $8$ khu nhỏ cho 8 đội. Diện tích mỗi khu nhỏ là bao nhiêu mét vuông?
  \shortans{300}
  \loigiai{
    Tổng diện tích là $60 \times 40 = 2\,400\text{ m}^2$. Mỗi khu có diện tích: $2\,400 : 8 = 300\text{ m}^2$.
  }
\end{ex}

% Câu 20
\begin{ex}
  Nhiệt độ ngoài trời lúc buổi sáng là $24^\circ\text{C}$, buổi trưa tăng thêm $8^\circ\text{C}$. Nhiệt độ buổi trưa là bao nhiêu độ C?
  \shortans{32}
  \loigiai{
    Nhiệt độ là: $24 + 8 = 32^\circ\text{C}$.
  }
\end{ex}

% Câu 21
\begin{ex}
  Một thước kẻ học sinh dài $20\text{ cm}$. Cần ghép bao nhiêu chiếc thước kẻ như thế nối tiếp nhau để được chiều dài $2\text{ m}$?
  \shortans{10}
  \loigiai{
    $2\text{ m} = 200\text{ cm}$. Số thước kẻ cần là: $200 : 20 = 10$ (chiếc).
  }
\end{ex}

% Câu 22
\begin{ex}
  Một cuộn giấy dán tường hình chữ nhật có kích thước $10\text{ m} \times 0{,}5\text{ m}$. Diện tích cuộn giấy đó là bao nhiêu mét vuông?
  \shortans{5}
  \loigiai{
    Diện tích cuộn giấy là: $10 \times 0{,}5 = 5\text{ m}^2$.
  }
\end{ex}

% Câu 23
\begin{ex}
  Người ta dán một bức tường diện tích $20\text{ m}^2$ bằng các cuộn giấy dán tường loại $5\text{ m}^2$/cuộn. Cần bao nhiêu cuộn giấy?
  \shortans{4}
  \loigiai{
    Số cuộn giấy cần là: $20 : 5 = 4$ (cuộn).
  }
\end{ex}

% Câu 24
\begin{ex}
  Một bể cá mini hình hộp chữ nhật có diện tích đáy là $12\text{ dm}^2$, chiều cao mực nước là $3\text{ dm}$. Thể tích nước trong bể là bao nhiêu lít (biết $1\text{ dm}^3 = 1\text{ l}$)?
  \shortans{36}
  \loigiai{
    Thể tích nước là: $12 \times 3 = 36\text{ dm}^3 = 36$ lít.
  }
\end{ex}

% Câu 25
\begin{ex}
  Thời gian một tiết học là $40\text{ phút}$, giờ ra chơi là $15\text{ phút}$. Hỏi thời gian $2$ tiết học và $1$ lần ra chơi là bao nhiêu phút?
  \shortans{95}
  \loigiai{
    Tổng thời gian là: $40 \times 2 + 15 = 80 + 15 = 95$ phút.
  }
\end{ex}
